\documentclass{situationist-preprint}

\title{Russellian Monism Without Panpsychism}
\msid{MND-18-b}
\msfamily{MND-18}
\msdiscipline{Philosophy of Mind}
\mstype{Companion manuscript}
\mskeywords{Russellian monism; panprotopsychism; protophenomenal properties; combination problem; intrinsic natures}

\begin{document}
\maketitle

\begin{abstract}
MND-18-a argues that panpsychism is false because the combination problem cannot be solved, and that the Russellian monism of MND-01 must therefore take the intrinsic natures of fundamental entities to be protophenomenal. We argue that the resulting view is stable. It is not a disguised physicalism, since protophenomenal natures are not structural; it is not a disguised panpsychism, since they are neither experiences nor subjects; and the combination problem it faces is a problem of articulating a grounding relation, not a proof that no such relation exists. The view retains what made Russellian monism attractive: a place for consciousness in the intrinsic nature of physical structure, and an explanation of the explanatory gap as a gap between modes of presentation.
\end{abstract}

\section{The position}

Russellian monism takes physics to describe the structure and dynamics of the world, its dispositional and relational profile, and leaves open the intrinsic natures that occupy that structure \citep{Russell1927}. MND-01 holds that consciousness is the intrinsic nature of physical structure, and that the explanatory gap arises because the same facts are given under two modes of presentation, structural and intrinsic. MND-18-a shows that this cannot be panpsychist: if the intrinsic natures of fundamental entities were their experiences, Theorem~3.2 of MND-18-a would rule out the grounding of our experience in them, and Corollary~3.3 draws the consequence. Proposition~4.1 shows that grounding in protophenomenal natures (Definition~2.3) escapes the argument.

\begin{definition}[Protophenomenal Russellian monism]\label{d:prm}
\emph{Protophenomenal Russellian monism} holds that (i) fundamental physical properties are structural and have intrinsic categorical bases; (ii) these bases are protophenomenal: not experiences, not subject-making, but jointly sufficient, in suitable structural arrangements, to ground experiences; and (iii) the experiences of macro-subjects are grounded in them \citep{Chalmers2013}.
\end{definition}

On this view the thesis of MND-01 is true at the level at which consciousness occurs: our experience is the intrinsic nature of a certain physical structure, the one our brains realize, and that intrinsic nature is grounded in the protophenomenal natures of the structure's fundamental constituents.

\section{Stability}

\begin{objection}[Collapse into physicalism]
Protophenomenal properties are characterized only by their grounding role. A property characterized by its role is structural, so the view is a form of physicalism and inherits its explanatory gap.
\end{objection}

\begin{reply}
The objection confuses how a property is introduced with what it is. Protophenomenal properties are introduced by a role, as the categorical grounds of physical dispositions and of experience, but the role is played by intrinsic properties, not by further structure: they are what the structure is a structure of. The distinction is the one between the two conceptions of the physical drawn by \citet{Stoljar2001}, and the view is physicalist only on the conception that includes the categorical grounds of physical dispositions. The explanatory gap does not return. On MND-01, we know the intrinsic nature of one physical structure under its intrinsic mode of presentation, in having our experiences, and the gap is the difference between that mode and the structural one. The protophenomenal grounds are what that intrinsic nature is grounded in; they do not need to be known under the intrinsic mode for the gap to be explained.
\end{reply}

\begin{objection}[Collapse into panpsychism]
If protophenomenal properties ground experience, they must be experiential in some attenuated way, and then the combination problem returns \citep{Goff2017}.
\end{objection}

\begin{reply}
Grounding does not transmit kind. Chemical properties ground life without being alive, and nothing in the grounding of experience requires its grounds to be experiences. What the objection can show is that protophenomenal grounding is not intelligible as composition of experience from experience. That is true, and it is the source of its immunity: the argument of MND-18-a turns on the fact that panpsychist grounds have owners (Principles P1--P3, Lemma~3.1), and protophenomenal grounds have none.
\end{reply}

\begin{objection}[The combination problem in another form]
The protophenomenal monist must still say which arrangements ground a subject, why one subject rather than many, and how a rich palette of qualities arises from a few fundamental natures \citep{Chalmers2017,Coleman2014}.
\end{objection}

\begin{reply}
She must, and these are the boundary and palette problems. But they are problems of spelling out the principles of a grounding relation, of the kind every theory of grounding faces, and none of them yields a contradiction with the principles that refute panpsychism (Proposition~4.1 of MND-18-a). The difference is that between a theory that has not yet been completed and one that cannot be completed. The treatment of the boundary and palette problems belongs to MND-01, whose companion MND-01-d gives a preliminary account of them.
\end{reply}

\section{What is kept}

Protophenomenal monism keeps three advantages of Russellian monism. It assigns consciousness a causal role, since the intrinsic natures are the categorical bases of the dispositions physics describes \citep{Chalmers2013}. It avoids the conceivability argument against physicalism, since zombie worlds duplicate structure without intrinsic natures. And it locates the explanatory gap in two modes of presentation rather than in two kinds of fact, as MND-01 requires. What it gives up is only the thesis that matter is everywhere conscious, which was never needed.

\section{Conclusion}

Panpsychism is false, and the Russellian monism of MND-01 survives without it. The intrinsic natures of fundamental entities are protophenomenal; our experience is the intrinsic nature of the structures they compose; and the combination problem, in the form that remains, is a task for the theory of grounding, not a refutation.

\begin{thebibliography}{9}\small
\bibitem[Chalmers(2013)]{Chalmers2013} Chalmers, D.~J. (2013). Panpsychism and panprotopsychism. \emph{The Amherst Lecture in Philosophy}, 8.
\bibitem[Chalmers(2017)]{Chalmers2017} Chalmers, D.~J. (2017). The combination problem for panpsychism. In G.~Br\"untrup and L.~Jaskolla (eds.), \emph{Panpsychism: Contemporary Perspectives}. Oxford University Press.
\bibitem[Coleman(2014)]{Coleman2014} Coleman, S. (2014). The real combination problem: panpsychism, micro-subjects, and emergence. \emph{Erkenntnis}, 79(1), 19--44.
\bibitem[Goff(2017)]{Goff2017} Goff, P. (2017). \emph{Consciousness and Fundamental Reality}. Oxford University Press.
\bibitem[Russell(1927)]{Russell1927} Russell, B. (1927). \emph{The Analysis of Matter}. Kegan Paul.
\bibitem[Stoljar(2001)]{Stoljar2001} Stoljar, D. (2001). Two conceptions of the physical. \emph{Philosophy and Phenomenological Research}, 62(2), 253--281.
\end{thebibliography}

\end{document}
